\chapterimage{apps.jpg}

\chapter{Capa 7: Aplicación}\label{sec:layer7}

\begin{minipage}{0.4\linewidth}
\begin{center}
\begin{bytefield}{16}
\bitbox{16}{\color{color1} Layer 7: Application} \\
\bitbox{16}{Layer 4: Transport} \\
\bitbox{16}{Layer 3: Internet} \\
\bitbox{16}{Layer 2: Network (LAN)} \\
\bitbox{16}{Layer 1: Physical} \\
\end{bytefield}
\end{center}
\end{minipage}
\begin{minipage}{0.6\linewidth}
\begin{center}
\includegraphics[width=\linewidth]{network_layer7.pdf}
\end{center}
\end{minipage}

La capa superior de la pila de red TCP/IP incluye todas las aplicaciones
que usan TCP o UDP como transporte. Es, de largo, la que contiene
el mayor número de protocolos, y donde esos protocolos se crean y se
reemplazan más rápidamente.

Además de las aplicaciones de usuario, existen varios sistemas y protocolos auxiliares que
ayudan a configurar y operar todos los hosts (equipos conectados a la red) de Internet.
En este capítulo se abordan algunos de los más críticos.

\section{DNS -- Domain Name System}\label{sec:layer7:dns}

El protocolo \concept{DNS} se encarga de traducir \conceptRef{domain name}{nombres de dominio}
como \otherBase{ietf.org} o\\\otherBase{cv.uab.cat} en direcciones \concept{IP}
que se pueden colocar dentro de \conceptRef{datagram}{datagramas}.

Los servicios DNS funcionan tanto sobre \concept{TCP} como sobre \concept{UDP}, \ie, se pueden habilitar
ambos protocolos de transporte. En ambos casos se usa el puerto \otherBase{53}.

\subsection{Jerarquía}\label{sec:layer6:dns:hierarchy}

\begin{center}
\includegraphics[width=0.65\linewidth]{dns_hierarchy.pdf}
\end{center}

Los nombres de dominio de Internet son secuencias de nombres separados por puntos.
Estos nombres se presentan de lo específico a lo general, de forma jerárquica.
Por ejemplo, en \otherBase{cv.uab.cat}, \otherBase{cv} pertenece a \otherBase{uab.cat},
y \otherBase{uab} a \otherBase{cat}.
% 
El nombre situado más a la derecha, \eg\ \otherBase{cat} o \otherBase{com},
debe ser uno de los \conceptRef{top-level domain}{dominios de primer nivel} de la IANA.

Los servidores DNS se organizan de la misma forma jerárquica que esos nombres de dominio.
Hay $13$ servidores \concept{root DNS} (\inlineCode{Root-A} a \inlineCode{Root-M}),
que son responsables de todos los dominios de primer nivel.
Puedes preguntar a cualquiera de ellos por \otherBase{com},
pero también por \otherBase{su}, y sabrán de ellos.

Eso no significa que un servidor raíz conozca todos los dominios que hay \textit{por debajo}
de su dominio de primer nivel (\eg\ \otherBase{uab.cat}): solo conocen esos
\href{https://en.wikipedia.org/wiki/List_of_Internet_top-level_domains}{\underline{dominios de primer nivel}}.
(\eg\ \otherBase{cat}).
% 
El trabajo de los servidores DNS raíz es redirigirte a otro servidor DNS encargado
del \conceptRef{top-level domain}{dominio de primer nivel} que has solicitado.
 
De forma similar, los servidores DNS responsables de \otherBase{cat} conocerán
cualquier dominio \otherBase{*.cat} existente (\eg\ \otherBase{uab.cat}),
pero no sus subdominios (\eg\ \otherBase{cv.uab.cat}).
% 
El trabajo de esos servidores DNS de dominio de primer nivel es redirigirte a
otros servidores DNS \conceptRef{authoritative}{autoritativos} para el dominio de segundo nivel
solicitado (\eg\ \otherBase{uab.cat}).
% 
Los dominios de segundo nivel pueden entonces estructurar su \conceptRef{zone}{zona} de nombres
como deseen, incluyendo cualquier número y profundidad de subdominios, así como
la traducción de nombre a IP \textit{dentro} de su zona de autoridad.

\begin{exercise}
Todo el mundo usa DNS constantemente. ¿Cuáles serían los riesgos de tener
un sistema centralizado de traducción de nombres en lugar de la naturaleza distribuida del DNS?
\end{exercise}

\subsection{Registros}\label{sec:layer7:dns:records}

El DNS puede verse como una base de datos distribuida con distintos tipos de entradas
llamadas \conceptRef{DNS record}{registros}. Los principales son:

\begin{itemize}
\item[\textbf{A}:] Una entrada \concept{A} asocia un nombre de dominio a una dirección IPv4.
\item[\textbf{AAAA}:] Una entrada \concept{AAAA} asocia un nombre de dominio a una dirección IPv6.
\item[\textbf{NS}:] Un \concept{NS} apunta a un servidor DNS que puede ayudar.
\item[\textbf{MX}:] Una entrada \concept{MX} proporciona el nombre de dominio del servidor de correo
  que gestiona los mensajes de un dominio.
\item[\textbf{CNAME}:] Una entrada \concept{CNAME} asocia un nombre de dominio a otro nombre de dominio.
\end{itemize}

Las aplicaciones hacen \conceptRef{request}{peticiones} a los servidores DNS indicando el tipo
de registros que les interesan. Las \conceptRef{reply}{respuestas} del servidor pueden ser para
los registros \concept{A} o \concept{MX} solicitados,
pero el cliente también puede recibir una redirección a otro DNS (\concept{NS})
o un alias a otro nombre de dominio (\concept{CNAME}).

\subsection{Resolución}\label{sec:layer7:dns:resolution}

La naturaleza distribuida y jerárquica de la base de datos DNS hace que
a veces sea necesario intercambiar varios mensajes para satisfacer una \conceptRef{query}{consulta} DNS.
 
Por ejemplo, Alice puede realizar una resolución completamente \conceptRef{iterative}{iterativa} por su cuenta.
Las dos primeras respuestas contienen una entrada \concept{NS} que apunta al siguiente
servidor DNS y la dirección IP de ese servidor DNS, pero no la dirección IP del
\otherBase{cv.uab.cat} solicitado.

\begin{center}
\includegraphics[width=\linewidth]{dns_iterative.eps}
\end{center}

El proveedor de Internet (\concept{ISP}) de Alice o su empresa ofrecerán probablemente servidores DNS
cuya tarea es recibir las peticiones de Alice y resolverlas por ella de forma \conceptRef{recursive}{recursiva}.

\begin{center}
\includegraphics[width=1.05\linewidth]{dns_recursive.eps}
\end{center}

La combinación de la resolución \conceptRef{recursive}{recursiva} y las tablas de \conceptRef{cache}{caché} hace
del DNS un sistema bastante eficiente, con una disponibilidad extremadamente alta y baja latencia.

\begin{exercise}
Puedes usar las herramientas \inlineCode{host} y \inlineCode{dig} para hacer consultas DNS
y ver los resultados. Por ejemplo, mi salida de \inlineCode{dig uab.cat -t ns}
contiene las siguientes líneas:

\vspace{0.25cm}
\begin{verbatim}
;; QUESTION SECTION:
;uab.cat.                       IN      NS
;; ANSWER SECTION:
uab.cat.                172800  IN      NS      dns.uab.es.
;; ADDITIONAL SECTION:
dns.uab.es.             172112  IN      A       158.109.0.1
\end{verbatim}
\vspace{0.25cm}

\begin{itemize}
\item ¿Qué tipo de registro solicité?
\item ¿Qué dos tipos de registro recibí?
\item ¿Por qué recibí dos tipos de registro?
\item Realiza una resolución iterativa manual de \otherBase{cv.uab.cat}.
  Puedes empezar con algo como \inlineCode{dig @202.12.27.33 cv.uab.cat}
  para consultar al servidor raíz M.
\end{itemize}
\end{exercise}


\section{DHCP -- Dynamic Host Configuration Protocol}\label{sec:layer7:dhcp}

Para participar plenamente en Internet, un host necesita tener
correctamente configurados los siguientes aspectos:
\begin{itemize}
\item \textbf{Dirección IP}:
La dirección IP es necesaria para identificar un nodo en la red, y para que ese nodo sepa qué mensajes van dirigidos a él.
\item \textbf{Máscara de red}:
Los nodos necesitan conocerla para determinar cuándo entregar los datagramas directamente y cuándo indirectamente.
\item \textbf{IP del gateway (pasarela)}:
Se necesita la dirección IP de al menos un router (encaminador) para entregar indirectamente mensajes con destino fuera de la LAN (red de área local) actual.
\item \textbf{Dirección IP del servidor DNS}:
Aunque no es estrictamente necesaria para usar TCP/IP, la dirección de un servidor DNS es extremadamente útil, e imprescindible para resolver
nombres de dominio en direcciones IP.
\end{itemize}

Tradicionalmente, estos parámetros había que configurarlos manualmente en cada ordenador.
Hoy en día, la mayoría de los ordenadores de usuario se configuran automáticamente con \concept{DHCP}.
Los sistemas operativos modernos incluyen un cliente DHCP que consulta periódicamente la LAN hasta quedar configurado.

\begin{remark}
El hardware de red viene con una dirección MAC por defecto preconfigurada.
Además, un nodo necesita conocer su dirección MAC para utilizar el protocolo DHCP.
Por tanto, las direcciones MAC no se configuran mediante DHCP.
\end{remark}

\newpage
DHCP funciona sobre UDP. Los servidores DHCP escuchan en el puerto \otherBase{67}, y los clientes
escuchan en el puerto \otherBase{68}. El intercambio DHCP ideal consta de $4$ mensajes:
Discovery, Offer, Request, Acknowledgement.

\begin{center}
\includegraphics[width=\linewidth]{dhcp_dora.eps}
\end{center}

\begin{itemize}

\item \textbf{Discovery}\\
El objetivo de este mensaje es encontrar servidores DHCP dispuestos a proporcionar
una configuración.\\[-0.3cm]

En la primera comunicación, el cliente
no conoce la IP del servidor DHCP. En su lugar, usa:
\begin{itemize}
\item Su dirección MAC como dirección MAC de origen.
\item El broadcast (difusión) de la LAN (\otherBase{ff:ff:ff:ff:ff:ff}) como dirección MAC de destino.
\item \otherBase{0.0.0.0} como IP de origen, lo que significa ``desconocida''.
\item La dirección de \concept{local broadcast} , \otherBase{255.255.255.255}, como IP de destino.
\end{itemize}

\vspace{0.25cm}
\item \textbf{Offer}\\
Cero o más servidores responderán al cliente con ofertas de configuración.
Estas configuraciones (\eg\ direcciones IP) no son propiedad del cliente
hasta que el servidor envía un Acknowledgement DHCP.\\[-0.3cm]

El servidor lleva la cuenta de un \concept{pool} (conjunto) de direcciones que asigna
dinámicamente a los clientes.
% 
Las direcciones IP se ofrecen durante un periodo de tiempo fijo, tras el cual expiran
a menos que se \conceptRef{renew}{renueven} (véase más abajo). Las direcciones que no se renuevan vuelven
al pool y se ofrecen de nuevo a los clientes cuando hacen peticiones.
\\[-0.3cm]

\begin{itemize}
\item El servidor puede usar su dirección IP real y la ofrecida al cliente.
\item El cliente sabe que es el destinatario porque el servidor se fija en la
dirección MAC de origen del Discovery del cliente.
\end{itemize}

\vspace{0.25cm}
\item \textbf{Request}\\
El cliente elige una de las ofertas recibidas y pide que se complete la
asignación.\\[-0.3cm]

\begin{itemize}
\item El cliente usa la dirección de broadcast local, \otherBase{255.255.255.255}, como \mbox{destino}.
\item El cliente mantiene \otherBase{0.0.0.0} como origen porque la dirección ofrecida todavía no es suya.
\end{itemize}

\vspace{0.25cm}
\item \textbf{Acknowledgement (ACK)}\\
El servidor confirma la petición. Una vez el cliente recibe este ACK,
empieza a usar la configuración ofrecida como propia.
Este ACK no debe confundirse con el de TCP.
\\[-0.3cm]

\vspace{0.25cm}
\item \textbf{Negative Acknowledgement (NACK)}\\
También es posible que el servidor deniegue la petición, \eg\ si
ha pasado demasiado tiempo desde la oferta. En este caso,
se envía en su lugar un Acknowledgement negativo (NACK). Si se recibe un NACK,
el cliente debe volver a la fase Discover.
\end{itemize}

Cuando un servidor DHCP \conceptRef{lease}{cede} una dirección a un cliente, lo hace
durante un periodo de tiempo limitado $T$. Este periodo lo fija el servidor DHCP y se puede configurar.
Una vez transcurrido $T/2$, los clientes pueden intentar renovar la configuración cedida. Para ello,
el ciclo Request-Acknowledgement se repite a lo largo del tiempo hasta que alguno de los dos extremos no logra completarlo a tiempo,
o el servidor DHCP envía un NACK.

\begin{center}
\includegraphics[width=\linewidth]{dhcp_renewal.eps}
\end{center}

\begin{exercise}
El típico router doméstico implementa tanto un cliente DHCP como un servidor DHCP.
Explica cómo es posible y por qué resulta útil.
\end{exercise}


\section{Sistemas autónomos}\label{sec:layer7:as}

El encaminamiento se realiza un \conceptRef{hop}{salto} a la vez, usando \conceptRef{routing table}{tablas de encaminamiento}
como se ha visto en la Sección~\ref{sec:layer3:routing}. Esas tablas deben coordinarse a gran
escala para garantizar la conectividad del grafo hasta todos los hosts de Internet.

\begin{center}
  \vspace{-0.2cm}
 \includegraphics[width=\linewidth]{as.pdf}
 \vspace{-0.75cm}
\end{center}


\subsection{Organización}\label{sec:layer7:as:organization}

Para facilitar esta coordinación, las LAN se agrupan en \conceptRef{AS}{sistemas autónomos} (\concept{AS}).
Estos gestionan bloques de direcciones IP no solapados (\otherBase{w.x.y.z/N}) cuya propiedad y gestión son independientes.
Muchas empresas, universidades e ISP tienen sus propios AS, que están
\href{https://2ip.io/analytics/asn-list/}{\underline{listados}} públicamente. Por ejemplo, la red de la UAB pertenece al
AS13041.

Los \conceptRef{IXP}{puntos de intercambio de Internet} (\concept{IXP}) son un tipo especial de AS cuyo propósito es interconectar otros AS,
en lugar de alojar servicios y/o clientes. Por supuesto, hacer funcionar estos IXP no es gratis, así que se establecen relaciones
comerciales con los AS interconectados.
%
Algunos IXP se comprometen a dar servicio a cualquiera
que cumpla sus requisitos básicos, no están controlados por ningún operador concreto y ofrecen a sus clientes precios transparentes y a menudo
equitativos. Este tipo de IXP se denomina \conceptRef{NP}{Punto Neutro} (\concept{NP}), en contraposición
a los puntos no neutros usados, \textit{e.g.}, solo por un único operador de red.

\newpage
Todos los AS operan \conceptRef{BR}{routers frontera} (\concept{BR}) para interconectar AS. Cuando un router frontera recibe un datagrama
con destino dentro del AS, lo reenvía a través de los \conceptRef{internal router}{routers internos} (IR) del AS.
De forma similar, cuando un host dentro de un AS necesita contactar con un host de otro AS, los datagramas se reenvían por routers frontera
hasta que llegan al AS de destino y, finalmente, a la LAN de destino.

\subsection{Configuración de la tabla de encaminamiento}\label{sec:layer7:as:routing_table_configurati}

Configurar manualmente las \conceptRef{routing table}{tablas de encaminamiento} de todos los routers frontera y routers internos
sería demasiado costoso y propenso a errores. En su lugar, se emplean distintos algoritmos para configurarlas automáticamente.
\begin{itemize}
\item Dentro del AS, los routers internos y los routers frontera usan uno de los
\conceptRef{interior gateway protocol}{protocolos de gateway interior} disponibles.
A menudo se usa \concept{RIP} (Routing Information Protocol) por su configuración relativamente sencilla.
Dicho esto, los AS pueden elegir libremente el algoritmo que usan internamente en su AS.\\[-0.2cm]

RIP implementa una variación del algoritmo de camino más corto de Bellman-Ford para minimizar el número de \conceptRef{hop}{saltos}
que los datagramas deben recorrer dentro del AS. Periódicamente (\eg\ cada $30$~s), cada IR y BR:\\[-0.4cm]
% 
  \begin{enumerate}
  \item Anuncia su distancia \textit{percibida} al resto de routers del AS.
  \item Actualiza sus distancias percibidas a partir de los anuncios de otros routers.
  \item Actualiza en consecuencia su \conceptRef{routing table}{tabla de encaminamiento} hacia todas las LAN del AS.
  \end{enumerate}

\begin{exercise}
RIP usa el número de saltos como métrica de coste.
% 
¿Por qué puede resultar problemático?
% 
¿Cómo se puede resolver este problema?
\end{exercise}

\vspace{0.25cm}
\item Entre AS, solo hay que configurar los routers frontera; los routers interiores simplemente necesitan dirigir
el tráfico exterior a través de los routers frontera. \\[-0.3cm]

Los routers frontera ejecutan un \conceptRef{exterior gateway protocol}{protocolo de gateway exterior} además del
protocolo de gateway interior usado dentro del AS. Los distintos AS deben ponerse de acuerdo en el
algoritmo que usan, así que normalmente se usa \concept{BGP} (border gateway protocol).\\[-0.3cm]

\concept{BGP} es un protocolo más complejo que \concept{RIP}. Esta complejidad
es necesaria porque BGP coordina routers frontera pertenecientes a distintos AS, \ie,
con distintos propietarios:\\[-0.4cm]
\begin{itemize}
\item RIP usa un único número (la distancia, \ie, el número de \conceptRef{hop}{saltos})
para guiar un algoritmo de camino más corto en el grafo. BGP usa \conceptRef{path}{caminos} en lugar de distancias.
Esto permite a BGP elegir rutas con más \conceptRef{hop}{saltos}, pero más deseables para el AS,
\eg\ en términos de tiempo, coste económico, carga interna, etc.\\[-0.4cm]

\item BGP se diseñó teniendo en mente la seguridad y la confianza, e incluye mecanismos
para verificar la autenticidad de los caminos anunciados por un router frontera.
\end{itemize}
\end{itemize}

\begin{remark}
BGP también puede usarse como protocolo de gateway interior del AS.
Normalmente esto solo lo hacen AS muy grandes.
\end{remark}

\begin{exercise}
¿Cómo puede BGP traducir el interés económico en configuraciones de encaminamiento?
\end{exercise}

